\solution{Магниттік тізбектер және трансформатор}{10.0}

\tasksol Бұл --- шексіз ұзын соленоидтың магнит өрісі туралы стандартты есеп. Ұзындығы $l$ болатын тікбұрышты контур үшін Ампер теоремасы бойынша

\begin{eq}[points=0.2]
    H\cdot l = NI,
\end{eq}
және өзекше ішіндегі $B=\mu\mu_0 H$ ескерсек, аламыз

\begin{eq}[points=0.1, numpoints=0.1]
B = \frac{\mu\mu_0 NI}{l}=\qty{1}{\T}.
\end{eq}

\criterion[type=remark]{Егер қатысушы \eqref{1}-ді бірден $B=\mu\mu_0 H$ арқылы жазса, \eqref{1} үшін де балл беріледі}

\tasksol Магнит ағыны келесідей берілген:

\begin{eq}[points=0.2]
    \Phi = B\cdot \pi r^2.
\end{eq}

\eqref{2}-мен біріктіре отырып, аламыз

\begin{eq}[points=0.2, numpoints=0.2]
    R_m = \frac l{\mu\mu_0\cdot\pi r^2}=\qty{6.33e5}{\A\per\Wb}.
\end{eq}

Мұнда $\unit{\A\per\Wb}=\!\unit{\A\per\T\per\m\squared}=\unit{\per\H}$ бірліктері эквивалентті.

\tasksol Индуктивтілік жіберілген ток бірлігіне шаққандағы $\Psi$ толық ағынына (потокосцепление) тең. Мұнда

\begin{eq}[points=0.2]
    \Psi=N\Phi,
\end{eq}

сондықтан

\begin{eq}[points=0.2, numpoints=0.2]
    L = \frac{N^2}{R_m} = \frac{\mu\mu_0\pi r^2N^2}l =\qty{63.2}{\mH}.
\end{eq}

\tasksol Шектік шарттарды қарастыру үшін екі фактіні пайдалану керек: магнит өрісінің соленоидтылығы $\nabla\cdot\vec B=0$, бұл екі ортаға да қарап тұрған шағын цилиндрге Гаусс теоремасын қолданғанда \textit{$\vec B$ векторының нормаль құраушысының сақталуына} әкеледі; сондай-ақ бөлу шекарасына перпендикуляр қиып өтетін тікбұрышты контур үшін $\oint\vec H\cdot d\vec l=0$ қолданғанда $\vec H=\vec B/\mu\mu_0$ \textit{векторының тангенциал құраушысының сақталуы} орын алады. $\vec B$ векторының нормаль құраушысының сақталуы мынаны береді:

\begin{eq}[points=0.3]
    B_0\cos\alpha=B\cos\beta,
\end{eq}

ал $\vec H$ векторының тангенциал құраушысының сақталуы (ауа үшін $\mu=1$ болғандықтан):

\begin{eq}[points=0.3]
    B_0\sin\alpha = \frac B\mu \sin\beta.
\end{eq}

\eqref{7}--\eqref{8} теңдеулерін бірге шеше отырып, аламыз

\begin{eq}[points=0.2]
    \tan\beta = \mu \tan\alpha,
\end{eq}

яғни

\begin{eq}[points=0.1, prefix={Сандық мәні}]
    \beta \approx \ang{89.83},
\end{eq}

сонымен қатар

\begin{eq}[points=0.2, numpoints=0.1]
    \frac B{B_0}=\sqrt{\cos^2\alpha+\mu^2\sin^2\alpha} \approx 347.
\end{eq}

Көріп отырғанымыздай, радиандық өлшемдегі $\alpha\ll1$ жағдайынан басқа кез келген $\alpha$ бұрыштары үшін $\beta$ әрқашан 90 градусқа өте жақын, ал $B_0$ мәні $B$-ден айтарлықтай кіші. Бұл мұндай есептерде ағынның сыртқа шығып кетуін (утечка) ескермеуге болатынын көрсетеді.

\tasksol Ампер теоремасын өзекшенің ортасы бойымен қолданамыз:

\begin{eq}[points=0.2]
    H_1\cdot(4a-d) + H_2\cdot d = NI,
\end{eq}

мұндағы $H_1$ --- өзекше ішіндегі магнит өрісінің кернеулігі, ал $H_2$ --- ауа саңылауындағы мәні. Басқаша айтқанда, $B$ сақталатындықтан,

\begin{eq}[points=0.3]
    B = \mu_0 H_2 = \frac{\mu_0}\mu H_1.
\end{eq}

Сонымен, ағын үшін келесіні аламыз:

\begin{eq}[points=0.3]
    \Phi = \frac{\mu_0 NI S}{d + \mu(4a-d)}.
\end{eq}

Алайда, егер \eqref{4} формуласымен салыстырсақ, келесілерді қолдануға болатынын көреміз:

\begin{eq}[points=0.1]
    R_{m,1} = \frac{4a-d}{\mu\mu_0S},
\end{eq}

\begin{eq}[points=0.1]
    R_{m,2} = \frac d{\mu_0S},
\end{eq}

бұл $R_m=R_{m,1}+R_{m,2}$ деп жазып, есеп шартындағы (1) формуланы қолданумен аналогты.

\tasksol $R_{m,1}=R_{m,2}$ деп алсақ,

\begin{eq}[points=0.2, numpoints=0.2]
    d = \frac{4a}{\mu+1}\approx\frac{4a}{\mu} = \qty{0.4}{\mm}.
\end{eq}

Бұл бөлім магниттік кедергінің негізінен ауа саңылауларында шоғырланатынын айқын көрсетеді!

\tasksol%
\criterion{\solref{7}, \solref{10}--\solref{11} пункттерінде магниттік тізбектер заңдарына қатысты критерийлер (\eqref{18}, \eqref{20}, \eqref{27}--\eqref{29}, \eqref{33}--\eqref{34}) ұқсас өрнектер болған жағдайда қабылданады}%
%
Алдымен екі катушканың тек геометриялық қасиеті болып табылатын өзара индукция коэффициентін анықтайық. Сондықтан қарапайымдылық үшін тек $V(t)$ арқылы туындайтын $I_1$ тогы бар, ал $I_2=0$ деп есептейік. Бірінші реттік орама тудыратын ағын

\begin{eq}[points=0.2]
    \Phi(t) = \frac{N_1 I_1(t)}{R_m},
\end{eq}

мұндағы $R_m$ мәні $d=0$ үшін \eqref{15}-тен анықталады. Бұл екінші реттік орамада келесідей толық ағынды тудырады

\begin{eq}[points=0.2]
    \Psi_2(t) = N_2\Phi(t),
\end{eq}

сонда өзара индукция коэффициентінің $M = \Psi_2/I_1$ анықтамасынан аламыз

\begin{eq}[points=0.2]
    M = \frac{N_1N_2}{R_m} = \frac{\mu\mu_0N_1N_2S}{4a}.
\end{eq}

Енді нақты тізбекке қарап, эквивалентті кедергіні анықтайық. Фарадей заңы бойынша

\begin{eq}[points=0.2]
    V(t) = -N_1\frac{d\Phi}{dt},
\end{eq}

ал $V_2=I_2R$ кернеуі дәл сол $V_2=-N_2\dot\Phi$ заңы арқылы туындайды. Бұдан келесіні көреміз:

\begin{eq}[points=0.2]
    I_2 R = \frac{N_1}{N_2}V(t).
\end{eq}

Ағын магниттік тізбек үшін Ом заңымен анықталады: бірінші және екінші реттік орамалар магнит қозғаушы күшке (МҚК) үлес қосады, бірақ тізбектің кедергісі нұсқауда айтылғандай өте аз, сондықтан

\begin{eq}[points=0.2]
    N_1I_1+N_2I_2 = 0.
\end{eq}

Осыдан $I_1$ ток күшін өрнектей аламыз. \eqref{21}--\eqref{23} теңдеулерін шеше отырып, аламыз

\begin{eq}[points=0.3, override={R'=R\ab(\frac{N_1}{N_2})^2}]
    I_1 =-\frac{V(t)}{R}\ab(\frac{N_2}{N_1})^2\implies R'=R\ab(\frac{N_1}{N_2})^2.
\end{eq}

Минус таңбасы ток күшінің $V(t)$ айнымалы кернеуіне қатысты фаза бойынша $\pi$-ге ығысқанын көрсетеді.

\tasksol Магниттік кедергі туралы анықтамадан көрініп тұрғандай, магнит ағыны --- МҚК-нің магниттік кедергіге қатынасы. Электр өрісінің кернеулігі Фарадей заңы бойынша электр қозғаушы күштің әсерінің нәтижесі $\oint\vec E\cdot d\vec l =\mathcal E$, бұл Ампер теоремасына $\oint\vec H\cdot d\vec l = NI$ ұқсас. Магниттік өтімділік өріс кернеулігі $\vec H$ мен магниттік индукция ағынының тығыздығы $\vec B$ арасында байланыс орнатады, яғни соңғысы $\vec j =\sigma \vec E$ қатысы арқылы электр өрісімен байланысқан ток тығыздығы $\vec j$-ге ұқсас. Сонымен біз келесіні аламыз:

\begin{center}
    \begin{tabular}{|c|c|}\hline
         \textbf{Магниттік тізбек}& \textbf{Электрлік тізбек} \\\hline
         Магнит ағыны $\Phi$ &Ток күші $I$\\\hline
         Магнит өрісінің кернеулігі $\vec H$&Электр өрісінің кернеулігі $\vec E$ \\\hline
         Абсолют магниттік өтімділік $\mu\mu_0$&Меншікті өткізгіштік $\sigma\equiv1/\rho$\\\hline
    \end{tabular}
\end{center}

\criterion[points=0.1]{Жазылған: <<Ток күші $I$>>}
\criterion[points=0.2]{Жазылған: <<Магнит өрісінің кернеулігі $\vec H$>>}
\criterion[points=0.2]{Жазылған: <<Меншікті өткізгіштік $\sigma$ немесе $1/\rho$>>}

\textit{Қазылар алқасына ескерту: бірінші кезекте қатысушы жазған шаманың сөзбен аталуына, содан кейін ғана оның әріптік белгіленуіне басымдық беріледі. Термин жазылмаған жағдайда, аналогия үшін балл тек әріптік жазба авторлық нұсқамен дәл сәйкес келсе ғана қойылады.} Айта кету керек, аналогиялар \textit{тек} электрлік/магниттік тізбектер контекстінде ғана жүреді. Сондықтан, мысалы, диэлектрлік өтімділік $\varepsilon$ магниттік өтімділік $\mu$-ге ұқсас деп айтуға болмайды, тіпті басқа контексте бұл аналогия дұрыс болса да.

\tasksol Бізде магниттік энергия тығыздығының формуласы бар:

\begin{eq}[points=0.2]
    w_m = \frac{BH}{2} = \frac{B^2}{2\mu\mu_0}=\frac{\mu\mu_0 H^2}{2}.
\end{eq}

Егер ұзындығы $l$ және көлденең қимасының ауданы $S$ болатын учаскені қарастырсақ, $W_m=w_m Sl$ болады, және \eqref{4}-ті және \solref{8}-дегі аналогияны ескере отырып, қорытамыз

\begin{eq}[points=0.3, override={W_M = \frac{\mathcal F^2}{2R_m}}]
    W_M = \frac{\mathcal F^2}{2R_m}=\frac12\Phi^2R_m.
\end{eq}

\pic{8}R{0pt}{0.3}{figures/3.7.png}

\tasksol Эквивалентті электр тізбегін сызып, фаза мен нөлдің МҚК-і сыртқы контурда сағат тілі бағытымен ағын тудыратынын көруге болады. Мұнда $R_m=\rho_m\cdot l$ --- ұзындығы $l$ учаскенің кедергісі (бүкіл фигураның жалпы ұзындығы $7l$). $I_1=I_2$ болғанда орталық ``көпіршеде'' ағын нөлге тең болар еді, ал нөлдік емес мәнде осы пердеше арқылы өтетін $\Phi_1-\Phi_2$ ағындарының айырмасы бар. Екінші реттік орама тек келетін ағынға ``жауап береді'', сондықтан эквивалентті тізбекте бұл учаске (екінші реттік орама, $M$ магниттік іске қосқышы және $R$ кедергісі) сол күйінде қалады және техникалық тұрғыдан магниттік тізбектің элементтері болып табылмайды. Енді біз, мысалы, сыртқы контур үшін Кирхгофтың екінші ережесін жаза аламыз:

\begin{eq}[points=0.3]
    NI_1(t) + NI_2(t) = 3R_m(\Phi_1+\Phi_2)
\end{eq}

және, мысалы, сол жақ контур үшін:

\begin{eq}[points=0.3]
    NI_1(t)+N_0I_M(t) = 3R_m\cdot\Phi_1 + R_m\cdot (\Phi_1-\Phi_2).
\end{eq}

\eqref{27}--\eqref{28} шешімінен $\Delta\Phi=\Phi_1-\Phi_2$ үшін алатынымыз

\begin{eq}[points=0.3]
    \Delta\Phi = \frac{N(I_1(t)-I_2(t))}{5R_m}.
\end{eq}

Магниттік іске қосқыштың импедансы жоқ деп айтылғандықтан, екінші реттік орамадағы барлық кернеу резисторға кетіп, $I_M$ тогын тудырады. Фарадей заңы арқылы

\begin{eq}[points=0.1]
    I_M(t) = -\frac{N_0}R\frac{d}{dt}\Delta\Phi,
\end{eq}

(минус таңбасы міндетті емес, өйткені бізге амплитуда маңызды), және $\frac{d}{dt}I(t) = 2\pi f\cdot I(t+\pi/2)$ екенін ескеріп, \eqref{29}-ға қолдансақ, амплитудалық мән үшін аламыз

\begin{eq}[points=0.3, numpoints=0.2]
    I_M = \frac{2\pi fNN_0(I_1-I_2)}{\sqrt{(5R\rho_m l)^2+(4\pi fN_0^2)^2}}=\qty{11.8}{\mA}.
\end{eq}

\tasksol $x$ --- $\mathcal F$ қосылған кездегі серіппелердің сығылуы болсын. Бұл жағдайда саңылауда келесі көлем қалады:

\begin{eq}[points=0.1]
    V=\pi a^2(d-x).
\end{eq}

$\mu\gg1$ болғандықтан, \solref{9}-дан магниттік энергия бекітілген ағын үшін $R_m$-ге пропорционал екенін көреміз, ал \solref{6}-дан магниттік кедергі негізінен ауа саңылауы арқылы туындайтынын ескерсек, магниттік энергияның көп бөлігі дәл осы ауа саңылауында жинақталады деп жоғары дәлдікпен айта аламыз. \textbf{(b)} және \textbf{(c)} өзекшелері арасындағы кеңістік үшін

\begin{eq}[points=0.1]
    W_m = \frac{\mathcal F^2}{2R_m},
\end{eq}

мұндағы $R_m$ келесіге тең:

\begin{eq}[points=0.4]
    R_m = \frac{d-x}{\mu_0 \pi a^2}.
\end{eq}

\criterion[type=remark]{Қатысушы өзекшелерден немесе бүйірлік саңылаулардан келетін энергия үлестерін жаза алады; балл олардың өрнегінде дұрыс \eqref{33}--\eqref{34} қосылғышы болған жағдайда қойылады}

Виртуалды $x$ орын ауыстырулары энергия градиентін тудырады, бұл пондеромоторлық күш ретінде көрінеді:

\begin{eq}[points=0.2]
    F=-\frac{dW_m}{dx},
\end{eq}

немесе \eqref{33}-тен,
$$
F(x) = \frac{\mu_0\pi a^2\mathcal F^2}{2(d-x)^2},
$$
яғни $F$ күші $x$-тің азаю бағытында әсер етеді. Жалпы айтқанда, магнит өрісі тек ауа саңылауында ғана емес, кеңістіктің басқа бөліктерінде де бар, бірақ бұл аймақтардың геометриясы якордың қозғалысы кезінде өзгермейді. Демек, тиісті энергия $x$-ке тәуелді емес және $F$ күшіне әсер етпейді. Сонымен, тұйықталу шарты

\begin{eq}[points=0.2]
    F(x_c) \geq kx_c.
\end{eq}

Магниттік энергияның өзгеруін серіппе энергиясымен $kx^2/2=\Delta W_m$ деп теңестіру бұл жерде қолданылмайтынын ескеріңіз. Мұндай қатынас магниттік күш әсерінен еркін қозғалыс кезіндегі якордың максималды ауытқуын сипаттайды. Біздің есепте іске қосылу шартын, яғни магниттік күш серіппе күшінен кем болмайтын сәтті анықтау талап етіледі. Сондықтан олардың энергияларын емес, $F(x)$ және $kx$ күштерін салыстыру қажет. Сонымен, \eqref{33}-ті \eqref{35} бойынша дифференциалдап және \eqref{36}-ға қойып, аламыз

\begin{eq}[points=0.5, override={\mathcal F_0 = \frac{d-x_c}a\sqrt\frac{2kx_c}{\mu_0\pi}}]
    \mathcal F_0 = \frac{d-x_c}a\sqrt\frac{2kx_c}{\mu_0\pi}=\qty{15.75}{\A}.
\end{eq}

$\mathcal F_0=NI_M$ қолданып, сандық түрде келесіні аламыз:

\begin{eq}[points=0.2, prefix={Сандық мәні}]
    I_M = \qty{11.25}{\mA}.
\end{eq}

\eqref{31} мәні \eqref{38}-ден сәл асатыны көрініп тұр (дифференциалдық трансформатор тудыратын ток іске қосылу шегінен сәл жоғары), сондықтан берілген жағдайларда сақтандырғыш іске қосылуы керек.

\criterion[points=0.1]{Көрсетілген, $I_M<\qty{11.8}{\mA}$}
